This study evaluates the performance of a proportional-integral-derivative
(PID) control system for precise temperature regulation in an open system
using a cost-effective experimental setup.
The system comprises a thermistor for temperature sensing, a heating element,
and an Arduino microcontroller to implement the PID algorithm.
The PID approach dynamically adjusts the power supplied to the heating element
by accounting for proportional (\(k_p\)), integral (\(k_i\)), and derivative
(\(k_d\)) terms, enabling rapid and stable responses to temperature
fluctuations caused by environmental changes.
Experimental results demonstrated that the PID-controlled system maintained a
stable temperature with a standard deviation of \( \sigma = x \) and a percent
error of XX\%, significantly outperforming traditional on-off control methods.
Furthermore, the data were used to calculate the specific heat capacities of
aluminum (\(\frac{J}{g^{\circ}C}\)) and brass (\(\frac{J}{g^{\circ}C}\)),
showcasing the system's utility in thermodynamic measurements.
These findings underscore the potential of PID temperature control for
applications requiring high precision, adaptability, and energy efficiency in
dynamic environments.